Ecosystem modelling is an important tool for understanding and managing complex environments, with food web models being essential for representing trophic interactions and energy flows in marine ecosystems and predicting their responses to external pressures \citep{Christensen2004, Colleter2015,audzijonyte2019atlantis}. These models provide quantitative insights into ecosystem structure and function, enabling researchers to assess cumulative impacts of multiple stressors and support ecosystem-based management decisions \citep{Coll2015, Villasante2016,Geary2020}. However, constructing these models presents significant challenges, particularly in developing food webs that capture the complex web of trophic interactions within an ecosystem.

Traditional approaches to food web construction rely heavily on extensive literature review, data collation and expert knowledge, which are time-consuming and resource-intensive \citep{Holden2024a}. The process of assembling diet matrices is particularly challenging, requiring synthesis of diverse data sources including field studies, literature reviews, and expert opinion. This creates a significant bottleneck in ecosystem model development, especially when applying models to new geographical contexts \citep{Holden2024b}. Recent advances in artificial intelligence (AI) offer new opportunities to streamline the food web construction process and avoid such bottlenecks \citep{spilliasfuture}. AI tools have demonstrated success in both knowledge/evidence synthesis tasks \citep{spillias2024human,keck2025extracting,castro2024large,spillias2024evaluating,Zheng2023,nugraha2024traditional}, ecological and environmental tasks \citep{Fernandes2024,Li2024,Chen2024,dorm2025large,Noleto2024} and modelling tasks \citep{Lapeyrolerie2022, Tuia2022, Karniadakis2021}, but their application to process-based ecosystem modelling remains nascent. The key challenge lies in ensuring that AI-driven approaches can effectively synthesise available information while maintaining ecological validity.

We present a flexible framework for assembling and synthesizing user-defined and online resources to construct food webs using Large Language Models (LLMs). Our approach integrates multiple data sources, including global biodiversity databases, species interaction repositories, and locally-held unstructured or structured text, to automate key steps in food web development. The framework employs user-selected LLMs to group species into functional units and estimate trophic interaction strengths. We evaluate the system in four distinct Australian marine ecosystems - the Northern Australia, South East shelf, and South East offshore regions, where we assess the reproducibility of the approach, and in the Great Australian Bight where we assess the accuracy of the approach. Specifically, we test the precision (repeatability) and scientific accuracy of automated species grouping decisions, and the precision and accuracy of the resulting diet matrix proportions, with accuracy defined in terms of similarity to expert estimates. These regions offer contrasting environmental conditions, species assemblages, and ecological dynamics, providing a robust test of the framework's adaptability and reliability.

Our framework operates through a five-stage process: (i) species identification using the Ocean Biodiversity Information System (OBIS) to generate comprehensive species lists for defined regions; (ii) data harvesting from FishBase, SeaLifeBase, and Global Biotic Interactions (GLOBI) databases to collect ecological and trophic information; (iii) species grouping using LLMs to classify species into functional groups based on ecological roles; (iv) food web construction using retrieval-augmented generation to synthesize species-level data into group-level diet compositions; and (v) validation through multiple iterations to assess reproducibility and comparison with expert-derived models to evaluate accuracy. For reproducibility assessment, we conducted five independent iterations per region, calculating consistency scores and Jaccard similarity indices to measure grouping stability, and stability scores to evaluate trophic interaction consistency. For accuracy assessment, we compared our framework's outputs with an expert-derived Ecopath model of the Great Australian Bight ecosystem \citep{Fulton2018}, evaluating both taxonomic grouping decisions and diet matrix proportions. This methodological approach allows us to systematically evaluate both the precision and ecological validity of our AI-driven food web construction framework.
